\subsection{Local finite conditions and rank-four rigidity}
\label{sec:rh-symplectic-point-rigidity}

Retain the actual ordinary-RH rank-four datum of
Theorem~\ref{thm:rh-symplectic-arithmetic-data} and the notation of
Theorem~\ref{thm:rh-symplectic-local}. In particular \(e\geq2\),
\(\ell\) is odd, the distinguished local set is \(\{\ell\}\),
and the residual Frobenius separation (BU8) holds. All points below
satisfy the original \texttt{PtIn} condition, with the actual
\texttt{OnComp} at \(\ell\). We do not assume finiteness of the
weight-zero algebra in this section.

\begin{lemma}[Continuity and finite coefficients]
Every characteristic-zero point of the complete normalized global
ring, and every dual-number lift of such a point from that ring,
has its image over a finite coefficient extension and is continuous.
Consequently these assertions hold for the representations in
\texttt{EPt} and \texttt{PtIn} over the dual numbers.
\end{lemma}

\begin{proof}
Write \(B\) for a complete Noetherian local coefficient algebra,
with residue field \(k\), and \(f:B\to\overline E\) for an
\(O\)-algebra map. For \(x\in\mathfrak m_B\) with \(f(x)\ne0\),
suppose \(v_\ell(f(x))\leq0\). Then \(f(x)^{-1}\) is integral
over \(O\). If \(P\) is its monic polynomial over \(O\), of
degree \(d\), the unit \(x^dP(1/x)\) of \(B\) maps to zero,
a contradiction. Thus all elements of \(\mathfrak m_B\) map
to elements of strictly positive valuation. Each element of
\(B\) is the sum of a representative from \(O\) of its residue
and an element of \(\mathfrak m_B\), so \(|f(B)|\leq1\).
Finite generation of \(\mathfrak m_B\) gives continuity.

For completeness let \(\mathfrak p=\ker f\). The domain
\(B/\mathfrak p\) is integral over \(O\). Hence the only prime
containing \((\ell)+\mathfrak p^2\) is the maximal ideal.
The quotient \(B/\mathfrak p^2\) is therefore complete for
the \(\ell\)-adic topology and finite modulo \(\ell\).
Successive approximation gives its finiteness over \(O\).
A derivation \(D\) at \(f\) kills \(\mathfrak p^2\), so the
map \(f+\epsilon D\), with \(\epsilon^2=0\), has bounded image
in a finite-dimensional algebra over a finite extension of \(E\).
The same approximation proves continuity of this map
(as in Proposition~\ref{prop:selmer-tangent-comparison}).
All the original points under consideration factor through
the complete normalized quotient constructed in the preceding
section. This proves the assertion for the representations,
not just for their traces. No claim is made here about
points of an arbitrary noncompleted graded algebra.
\end{proof}

\begin{proposition}[The local finite conditions]
\label{prop:rh-symplectic-local-finite}
Every actual characteristic-zero point and every actual dual-number
lift is crystalline at \(\ell\) after restriction to
\(\mathbb Q_\ell(\mu_\ell)\). In particular every actual tangent
cocycle to a characteristic-zero point \(r\) belongs to
\[
 H^1_f(\mathbb Q,\mathfrak{sp}_4(\overline E)_{\operatorname{Ad}r}).
\]
Here cohomology over \(\overline E\) is taken by finite coefficient
extension.
\end{proposition}

\begin{proof}
Let \(\chi_i=(\tau,\Psi,\alpha,1)_i\) for \(0\le i\le3\).
Over \(L^\circ[1/\ell]\) the prescribed eigenvalues
\(\chi_i(s_0)\) are distinct. The polynomial projectors of
\texttt{Ord} can be divided by their nonzero constant
denominators there. The resulting idempotents \(P_i\)
have rank one, the ordered flag is invariant under the
decomposition group, and its \(i\)-th graded character
\(\kappa_i\) has inertia character \(\chi_i\).
This also holds over dual-number coefficient algebras:
the denominators are invertible scalars and the images
of the idempotents are free of rank one over the local
coefficient algebra.

One must check the labels of the Frobenius eigenvalues;
distinct weights on \(s_0\) alone do not prove the needed
nonresonance. The element \(\mathrm{ph}\) has, in the universal
local representation, finite order \(d_0\) prime to \(\ell\)
dividing \(|k^\times|\), as in the proof of (SM3).
Thus each
\[
 b_i=\operatorname{tr}(P_i r(\mathrm{ph}))
       \quad\hbox{satisfies}\quad b_i^{d_0}=1.
\]
The ring \(L^\circ[1/\ell]\) is a domain, and \(O\) contains the
Teichm\"uller roots of this order. Factoring the polynomial
\(X^{d_0}-1\) shows that \(b_i\) is a single constant root.
Evaluation at the diagonal augmentation fixes that root:
\[
 \kappa_i(\mathrm{ph})=\chi_i(\mathrm{ph}).
 \tag{SR1}
\]
This identity is valid on the whole selected component,
hence also on its dual-number points.

Put \(F_\ell=\mathbb Q_\ell(\mu_\ell)\).
On inertia of \(F_\ell\) the four graded characters
are the cyclotomic powers with exponents \(e+1,e,1,0\).
Indeed \(\xi,\psi\) are unramified at \(\ell\), and the
Teichm\"uller cyclotomic character is now trivial.
The element \(\mathrm{ph}\) belongs to \(G_{F_\ell}\)
and its cyclotomic value is one. For \(i<j\) the ratio
\(\kappa_i/\kappa_j\) has positive cyclotomic exponent.
When that exponent is greater than one it differs from
\(\varepsilon\) already on inertia. When the exponent
is one its value on \(\mathrm{ph}\), by (SR1), is
the corresponding value of \(\alpha\), or of \(\beta\)
if \(e=2\). In either case this value differs from one
by (BU8). In particular none of these ratios equals
\(\varepsilon\) on \(G_{F_\ell}\).

We spell out the extension argument to include the actual
dual-number condition. Work over a sufficiently large finite
coefficient field. The ordinary flag over dual numbers
induces a flag of vector spaces over this field whose graded
pieces have dimension two (dimension one for an ordinary
point). On \(G_{F_\ell}\) each piece is an unramified
representation times the indicated cyclotomic power, since
the exact inertia characters are fixed, even in the
dual-number lift. Such pieces are crystalline. Inductively
suppose the preceding subrepresentation \(U\) of the flag
is crystalline and let \(V\) be the next graded piece.
Then \(W=\operatorname{Hom}(V,U)\) is crystalline and has
strictly negative Hodge--Tate weights in the convention
where \(\varepsilon\) has weight \(-1\).
The characters of \(W^*(1)\) in its induced filtration
are nontrivial by the preceding ratio calculation.
Thus \(H^0(F_\ell,W^*(1))=0\).
The local Euler characteristic and the crystalline
Bloch--Kato dimension formula give
\[
 \dim H^1(F_\ell,W)=
 [F_\ell:\mathbb Q_\ell]\dim W+\dim H^0(F_\ell,W)
 =\dim H^1_f(F_\ell,W).
\]
Every such extension is consequently crystalline.
This proves the induction, including the case \(e=2\).

The vector-space self-extension associated with a
dual-number tangent now splits over \(B_{\mathrm{cris}}\)
on \(G_{F_\ell}\). Its cocycle therefore has finite
local class after restriction. Restriction-corestriction,
with characteristic-zero coefficients, descends the
finite class to \(G_{\mathbb Q_\ell}\), first in
\(\operatorname{End}(r)\). Differentiation of the exact
fixed-similitude equation places the cocycle in
\(\mathfrak{sp}_4\). The involution
\[
 \theta(X)=-J^{-1}X^tJ
\]
commutes with \(\operatorname{Ad}r\), and its \(+1\)
eigenspace is \(\mathfrak{sp}_4\). The projection
\((1+\theta)/2\) thus gives the same finite condition
in this direct summand. Away from \(\ell\), the original
inertia kernel is killed by the representation and the
dual-number lift. Their inertia images in the relevant
cocycle calculation factor through the original finite
quotient of inertia. Averaging over this quotient in
characteristic zero gives an unramified local class.
These are exactly the finite local conditions.
\end{proof}

\begin{lemma}[Cyclotomic irreducibility without an image-size hypothesis]
\label{lem:rh-symplectic-cyclotomic-irreducibility}
Let \(r\) be an absolutely irreducible actual characteristic-zero
point. Choose an imaginary quadratic field \(K\) with a ramified
prime outside the finite ramification set of the datum, including
\(\ell\). Then
\[
 r|_{G_{K(\mu_{\ell^\infty})}}
 \quad\hbox{is absolutely irreducible}.
 \tag{SR2}
\]
In particular \(r|_{G_{\mathbb Q(\mu_{\ell^\infty})}}\)
is absolutely irreducible.
\end{lemma}

\begin{proof}
We use the following elementary observation. If \(H\) is a normal
subgroup of \(G\) with abelian quotient and an absolutely
irreducible finite-dimensional characteristic-zero representation
\(V\) of \(G\) becomes reducible on \(H\), then \(V\) has a
nontrivial finite-order self-twist through \(G/H\).
Indeed the sum of the \(G\)-translates of a simple \(H\)-submodule
is all of \(V\), so restriction to \(H\) is semisimple.
The trace-zero subspace of \(\operatorname{End}_H(V)\)
is then nonzero and stable under \(G/H\) acting by conjugation.
Over an algebraically closed field the commuting linear
operators of this abelian action have a common eigenvector
\(T\ne0\) in that subspace. Its character \(\eta\) is nontrivial
by Schur's lemma. The relation
\[
 r(g)Tr(g)^{-1}=\eta(g)T
\]
makes \(T\) an intertwiner between irreducible representations,
hence invertible. Determinants give \(\eta^{\dim V}=1\).
The character is continuous when \(r\) is, by taking a nonzero
matrix coordinate of \(T\) in this equality.

Apply the observation with \(H=G_{K(\mu_{\ell^\infty})}\).
This is an abelian quotient of \(G_{\mathbb Q}\), and
\(K\) is disjoint from the cyclotomic extension since it
ramifies at a prime outside \(\ell\).
A self-twist is trivial on inertia at the chosen ramified
prime, since \(r\) is unramified there. Its quadratic
factor is therefore trivial, so the self-twist factors
through the cyclotomic extension alone.
On inertia at \(\ell\) a finite self-twist must preserve
the multiset of four graded characters of \(r\).
The four distinct cyclotomic exponents \(e+1,e,1,0\)
cannot be interchanged by a finite-order character:
restrict further to the kernel of all the finite
characters and use the infinite image of the cyclotomic
character. Each of the four inertia characters must
therefore be fixed individually by the self-twist.
It follows that the self-twist is trivial on inertia at
\(\ell\), which surjects onto
\(\operatorname{Gal}(\mathbb Q(\mu_{\ell^\infty})/\mathbb Q)\).
This contradicts the nontriviality of \(\eta\) and proves (SR2).
A suitable \(K\) always exists, for example
\(\mathbb Q(\sqrt{-q})\) for a prime \(q\equiv3\bmod4\)
outside the finite ramification set.
\end{proof}

\begin{theorem}[The rigidity implication for the remaining RH points]
\label{thm:rh-symplectic-rigidity}
Suppose an absolutely irreducible actual characteristic-zero
point \(r\) of the remaining ordinary-RH datum satisfies
the source automorphy condition \texttt{IsAut4Alt} (equivalently
the condition \texttt{IsAutQA} used there).
Then it satisfies the original \texttt{Raw.RigidAt} condition.
Thus the RH specialization of the rigidity implication
\texttt{RigiditySymplecticFour} holds, with its original
automorphy premise, without an additional residual adequacy
assumption.
\end{theorem}

\begin{proof}
Choose \(K\) as in Lemma~\ref{lem:rh-symplectic-cyclotomic-irreducibility}.
The source \texttt{IsAutQA} applies to every imaginary quadratic
field where the restriction of the continuous representation
is irreducible. By (SR2) it supplies a regular algebraic
cuspidal representation \(\pi\) on \(\mathrm{GL}_4/K\)
attached to \(r|_{G_K}\).
The fixed similitude \(\tau=\varepsilon^{e+1}\psi\xi\omega_T^{-1}\)
is algebraic and is defined on \(G_{\mathbb Q}\).
The alternating pairing identifies \(r^\vee\otimes\tau\)
with \(r\), and \(r\) extends across \(K/\mathbb Q\).
The corresponding conjugate duality of \(\pi\), with the
algebraic character of \(\mathbb Q\) corresponding to the
similitude in the attachment normalization, follows at all
unramified places and hence globally by strong multiplicity one.
Thus \(\pi\) is polarizable. The polarization on the attached
Galois representation is symmetric: if \(c\) is complex
conjugation, \(\tau(c)=-1\), \(r(c)^2=1\), and
\[
 P=r(c)J^{-1},\qquad
 P^t=P,\qquad
 r(cgc^{-1})=\tau(g)P r(g)^{-t}P^{-1}\quad(g\in G_K).
 \tag{SR3}
\]
This is the polarized adjoint convention, with negative transpose
on the nonidentity component. Equivalently one may twist the
attached representation in passing to the chosen standard
algebraic normalization; this does not change the adjoint
representation or (SR2).

By Chebotarev and semisimplicity the attached representation
is the actual restriction (with that normalization).
The polarization is unique up to scalar by absolute irreducibility.
Thorne's adjoint vanishing theorem
\cite[Theorem~6.1]{support-type-one-thorne} now gives
\[
 H^1_f(\mathbb Q,\operatorname{ad}_{\mathrm{pol}}(r|_{G_K}))=0:
\]
its cyclotomic irreducibility hypothesis is precisely (SR2).
On \(G_K\) this adjoint is the usual conjugation action on
\(\operatorname{End}(r)\), and on \(c\) it is
\[
 X\longmapsto-PX^tP^{-1}
   =\operatorname{Ad}(r(c))\theta(X).
\]
The involution \(\theta\) commutes with both parts of this action.
On its \(+1\) eigenspace the action is exactly the usual
\(\operatorname{Ad}(r)\) on \(\mathfrak{sp}_4\).
The direct-summand projection therefore yields
\(H^1_f(\mathbb Q,\mathfrak{sp}_4{}_{\operatorname{Ad}r})=0\).
By Proposition~\ref{prop:rh-symplectic-local-finite},
every original dual-number tangent is a coboundary in this
summand, with a primitive \(X\in\mathfrak{sp}_4\).
Writing the tangent matrix as \(z(g)r(g)\), the equality
\(z(g)=X-\operatorname{Ad}(r(g))X\) gives
\(z(g)r(g)=Xr(g)-r(g)X\), exactly \texttt{Raw.RigidAt}.
\end{proof}

This argument proves local finite conditions for all actual points
and the full rigidity implication at automorphic irreducible RH
points. It does not prove the source's \texttt{AllIrr} automorphy
premise or \texttt{Fin0} or \texttt{NoEnd}. Potential crystallinity
does not by itself prove the all-powers nonresonance condition of
\texttt{PotUnl}. No integral Selmer length bound follows from the
characteristic-zero vanishing alone.
